\subsection{治疗的循证阶梯与常见误区}

现代指南对早泄的处理已有较清晰的"阶梯"共识，核心是\textbf{按需、按型、按困扰程度}选择手段，而非一味"加药加量"：

\begin{enumerate}
	\item \textbf{一线——教育与行为疗法}：性知识教育、消除表现焦虑，以及挤压法、停-动法、阴道静止法等训练。终身性早泄单靠行为疗法常难完全达标，但仍应作为一切治疗的基础。
	\item \textbf{一线——按需用药}：达泊西汀（性交前 1--3 小时，30 mg 或 60 mg）或局部麻醉剂（利多卡因/丙胺卡因乳膏或喷雾）。起效快、可控性好，适合多数患者。
	\item \textbf{二线——每日服药与联合治疗}：每日服用的 SSRI（帕罗西汀、舍曲林等，多为超说明书使用）；合并勃起功能障碍时联用 PDE5 抑制剂。行为疗法与药物联合，通常优于任何单一手段。
	\item \textbf{三线——手术，须极其谨慎}：\textbf{选择性阴茎背神经切断术已被主流指南明确不推荐}——其宣称的"一劳永逸"缺乏可靠的长期证据，且可能造成永久性龟头感觉减退，甚至勃起功能障碍，医学界对其争议极大。切勿轻信广告，务必到正规医院充分知情后再做判断。
\end{enumerate}

\noindent\textbf{几个常见误区}：

\begin{itemize}
	\item \textbf{"时间越长越好"}：性满意度与时长并非线性关系，越是执着于"计时"，越会加重表现焦虑，反而使早泄更糟。
	\item \textbf{"买瓶延时喷剂就搞定"}：市售延时喷剂、湿巾多含利多卡因、苯佐卡因等局麻成分，浓度与剂量不明，用量过大会导致龟头麻木、影响伴侣感觉，甚至引起不射精或勃起不坚；来源不明的产品还可能含违禁成分。
	\item \textbf{"纯中药壮阳无毒副作用"}：宣称"根治早泄"的壮阳偏方多缺乏循证依据，长期服用可能扰乱内分泌、损害肝肾。
	\item \textbf{把"偶尔过快"当病}：偶发的、情境性的射精过快属于正常波动（自然变异性早泄），不必用药；把"勃起维持困难时的抢时间"误当成早泄（早泄样射精功能障碍）也很常见，其处理重点是信心与认知，而非加药。
	\item \textbf{"自己练就行，不必让伴侣知道"}：研究一致显示，伴侣的参与能显著提高行为疗法的效果与长期维持率。
\end{itemize}

\subsection{伴侣参与、心理与关系因素}

早泄常被当成"男方一个人的问题"，但它几乎总是发生在关系之中、也在关系之中得到缓解。伴侣的反应——是理解配合，还是抱怨嘲讽——往往直接影响治疗效果。

\begin{itemize}
	\item \textbf{重新定义"成功"}：把目标从"坚持多少分钟"改为"双方是否满意、是否愿意继续亲近"。充分的前戏、协商的节奏、合适的体位（如女上位、侧位），以及必要时借助手法或器具补充，都能显著提高伴侣的满意度。
	\item \textbf{共同参与练习}：挤压法、停-动法最好由伴侣协助完成；约定一个"暂停信号"，让"停下来"不再尴尬，而成为双方共同遵守的游戏规则。
	\item \textbf{处理表现焦虑与关系冲突}：越是怕失败越容易失败。可借助放松训练与正念；若存在关系积怨、抑郁或性创伤，应接受认知行为治疗或专业性治疗师的帮助。
	\item \textbf{协同处理合并问题}：合并勃起功能障碍、慢性前列腺炎、甲状腺功能亢进或药物副作用者，须先治疗原发病——相当一部分"早泄"会在原发病得到控制后自行改善。
\end{itemize}

\begin{tcolorbox}[colback=pink!5!white,colframe=red!50!black,title=贴心小叮咛]
早泄是疗效最好的男性性功能障碍之一。请到正规医院男科或泌尿外科就诊，先分清类型、再谈治疗；切忌因羞于启齿而自行网购"神药"，也不要因一次不如意就给自己贴上"早泄"的标签。
\end{tcolorbox}
